\documentclass[11pt]{article}

\usepackage[T1]{fontenc}
\usepackage[utf8]{inputenc}
\usepackage{lmodern}
\usepackage[margin=1in]{geometry}
\usepackage{microtype}
\usepackage{booktabs,tabularx,array,longtable}
\usepackage{enumitem}
\usepackage{xcolor}
\usepackage{listings}
\usepackage{hyperref}

\definecolor{codebg}{RGB}{247,247,247}
\definecolor{codeframe}{RGB}{205,205,205}
\definecolor{codecomment}{RGB}{0,105,75}
\definecolor{codekeyword}{RGB}{25,65,145}

\lstdefinestyle{base}{
  basicstyle=\ttfamily\footnotesize,
  columns=fullflexible,
  keepspaces=true,
  showstringspaces=false,
  breaklines=true,
  breakatwhitespace=false,
  frame=single,
  rulecolor=\color{codeframe},
  backgroundcolor=\color{codebg},
  numbers=left,
  numberstyle=\tiny\color{gray},
  numbersep=7pt,
  xleftmargin=1.5em,
  framexleftmargin=1em,
  aboveskip=0.8em,
  belowskip=0.8em
}

\lstdefinestyle{cpp}{
  style=base,
  language=C++,
  keywordstyle=\color{codekeyword}\bfseries,
  commentstyle=\color{codecomment},
  stringstyle=\color{purple!65!black}
}

\lstdefinestyle{data}{
  style=base,
  language={}
}

\hypersetup{
  colorlinks=true,
  linkcolor=blue!55!black,
  urlcolor=blue!65!black,
  citecolor=blue!55!black,
  pdftitle={Data Structures and Memory Ownership in ParaDiS},
  pdfsubject={Developer guide to storage, aliases, lifetimes, and MPI ownership}
}

\setlist{nosep,leftmargin=*}
\setlist[description]{style=nextline,leftmargin=0pt,labelindent=0pt,
                      labelwidth=0pt,labelsep=0pt,
                      font=\normalfont\bfseries}
\setcounter{tocdepth}{2}
\emergencystretch=2em
\newcommand{\code}[1]{\texorpdfstring{\nolinkurl{#1}}{#1}}
\newcommand{\srcfile}[1]{\texorpdfstring{\nolinkurl{#1}}{#1}}

\title{Data Structures and Memory Ownership in ParaDiS\\
       \large Developer Guide to Storage, Aliases, Lifetimes, and MPI State}
\author{ParaDiS-Extended Development Notes}
\date{August 2026}

\begin{document}

\maketitle

\begin{abstract}
ParaDiS is a C/C++ hybrid whose production data structures are managed mainly
by procedural allocation, intrusive queues, sparse pointer tables, and
phase-specific MPI buffers.  A C++ pointer type therefore does not, by itself,
say whether a caller owns an object.  This guide documents the concrete
storage model: the per-rank \code{Home_t} root, \code{Param_t}, pooled
\code{Node_t} objects and their arm arrays, implicit and materialized segment
records, spatial cells, remote and mirror domains, inclusions, FMM state, and
important scratch structures.  It traces their allocation, reuse,
invalidation, migration, and final release paths.

The guide is written against repository commit \code{95b2842}.  It describes
the implementation as it exists, including several sharp edges that extension
code must respect.  Source files remain authoritative after later changes.
\end{abstract}

\tableofcontents
\newpage

\section{The ownership model in one page}

\subsection{Three meanings of ``owner''}

ParaDiS uses the word \emph{owner} for three different relationships.  They
must not be conflated.

\begin{description}
\item[Storage owner]
The component responsible for releasing an allocation.  For example,
\code{home->nodeBlockQ} owns production \code{Node_t} storage, while
\code{home->nodeKeys} only indexes those nodes.

\item[Domain or logical owner]
The MPI domain that holds the authoritative state for an entity.  A native
node has a local \code{myTag.domainID}; a ghost node is stored in local memory
but its tag names a remote logical owner.

\item[Operation owner]
The rank or endpoint selected to perform a calculation or topology mutation
once.  Routines such as \code{NodeOwnsSeg()} and \code{DomainOwnsSeg()} choose
work responsibility; they do not decide who frees memory.
\end{description}

Whenever this guide says \emph{owns} without qualification, it means storage
ownership.

\subsection{Per-rank ownership tree}

\begin{center}
\fbox{\parbox{0.92\linewidth}{
\ttfamily
Home\_t (one per MPI rank)\par
\quad |-- owns Param\_t and most rank-global arrays\par
\quad |-- owns NodeBlock\_t chain\par
\quad |\qquad `-- owns actual Node\_t slots\par
\quad |\qquad\quad `-- each active/recycled node owns arm arrays\par
\quad |-- owns Cell\_t hash entries and their integer arrays\par
\quad |\qquad `-- cell nodeQ borrows pooled Node\_t objects\par
\quad |-- owns RemoteDomain\_t records and lookup-array storage\par
\quad |\qquad `-- lookup entries borrow pooled ghost Node\_t objects\par
\quad |-- owns flat Eshelby inclusion storage when enabled\par
\quad |-- owns decomposition, FMM, operation-list, and scratch state\par
\quad `-- temporarily owns subcycling and communication state
}}
\end{center}

The same pooled node can be visible simultaneously through a tag lookup, a
ghost traversal array, and a cell queue.  These are aliases, not additional
owners.  Freeing through any of them would corrupt the pool.

\subsection{Lifetime classes}

\begin{longtable}{@{}p{0.18\linewidth}p{0.30\linewidth}p{0.44\linewidth}@{}}
\toprule
Lifetime & Typical structures & Boundary that invalidates or releases them \\
\midrule
\endhead
Process/rank & \code{Home_t}, \code{Param_t}, node blocks, persistent cell
hash entries & \code{ReleaseMemory()} at normal shutdown \\
Decomposition epoch & active \code{cellList}, cell/domain intersections,
\code{RemoteDomain_t} records & \code{DLBfreeOld()} followed by decomposition
and communication-neighbor rebuild \\
Outer cycle & ghost population, collision grid, some force and inclusion
scratch & migration, ghost recycle, resort, and the producer's clear routine \\
Algorithm phase & segment lists/pairs, migration buffers, topology receive
buffers, mirror snapshots & caller cleanup after the relevant computation or
after MPI completion \\
Call/local scope & temporary arrays and copied value records & return path of
the allocating routine \\
\bottomrule
\end{longtable}

\section{Source map}

\begin{longtable}{@{}p{0.38\linewidth}p{0.56\linewidth}@{}}
\toprule
File & Ownership responsibility \\
\midrule
\endhead
\srcfile{include/Home.h} & Per-rank root and memory-wrapper macros. \\
\srcfile{include/Param.h}, \srcfile{include/Typedefs.h} & Configuration,
parser registry records, and shared value types. \\
\srcfile{include/Node.h}, \srcfile{src/Node.cc}, \srcfile{src/Util.cc} & Node
layout and the C++ and legacy arm-management families. \\
\srcfile{src/QueueOps.cc}, \srcfile{src/GetNewNativeNode.cc},
\srcfile{src/GetNewGhostNode.cc} & Node-block allocation, pool transitions,
and active lookup installation. \\
\srcfile{include/Segment.h}, \srcfile{include/SegmentPair.h},
\srcfile{src/Segment.cc} & Derived segment views and C++ list ownership. \\
\srcfile{include/Cell.h}, \srcfile{src/CellTable.cc},
\srcfile{src/SortNativeNodes.cc}, \srcfile{src/DLBfreeOld.cc} & Persistent
cell storage and decomposition-epoch views. \\
\srcfile{include/RemoteDomain.h}, \srcfile{src/InitRemoteDomains.cc},
\srcfile{src/CommSendGhosts.cc} & Remote-domain metadata, ghost lookups, and
message buffers. \\
\srcfile{src/Migrate.cc}, \srcfile{src/RemapInitialTags.cc} & Ownership
transfer, retagging, and temporary tag maps. \\
\srcfile{src/RemoteOps.cc}, \srcfile{src/CommSendRemesh.cc},
\srcfile{src/FixRemesh.cc} & Persistent operation-list storage and transferred
receive-buffer ownership. \\
\srcfile{src/FMComm.cc}, \srcfile{src/SubcyclingIntegrator.cc} & FMM hierarchy
and C++ cycle-scoped subcycling state. \\
\srcfile{src/ParadisFinish.cc} & Ordered final cleanup. \\
\bottomrule
\end{longtable}

For graph mutation invariants and distributed topology replay, see
\srcfile{docs/Topology_and_Remesh_code.tex}.  For parameter registration and
broadcast details, see \srcfile{docs/Control_File_to_ParaDiS_Simulation_code.tex}.

\section{Home\_t: the per-rank root}

\subsection{Construction}

\code{InitHome()} deliberately zero-initializes the root.  Most pointer fields
therefore begin as null and most counters begin at zero.

\lstinputlisting[style=cpp,firstline=14,lastline=24,firstnumber=14]{../src/InitHome.cc}

Initialization then attaches independently allocated subsystems to this
root.  The important groups in \code{Home_t} are:

\begin{itemize}
\item rank, cycle, cached domain-boundary, and timing values;
\item the owned \code{Param_t} and rank-zero-only parser registries;
\item node-block ownership, node queues, lookup arrays, and the recycled-tag
heap;
\item persistent cell hash buckets plus the decomposition-epoch cell list;
\item remote-domain tables, communication request/status arrays, and buffers;
\item the generic decomposition pointer and optional FMM hierarchy;
\item topology operation storage, collision grids, tag maps, reduction
buffers, material tables, and solver state.
\end{itemize}

\code{Home_t} is an aggregation root, not an RAII container.  Its members use
different allocation families and require subsystem-specific cleanup such as
\code{FreeDecomp()}, \code{FMFree()}, and \code{FreeOpList()}.

\subsection{Procedural teardown}

Normal termination runs output and solver cleanup, releases MPI groups,
finalizes MPI, and then calls \code{ReleaseMemory()}.  The cleanup order is
intentional:

\begin{enumerate}
\item \code{DLBfreeOld()} invalidates decomposition-epoch cell and remote
aliases.
\item \code{FreeCellTable()} and subsystem destructors release nested
allocations.
\item every pooled node's remaining arm arrays and lock are released before
its raw node block is freed;
\item lookup arrays, heaps, operation storage, decomposition data, material
tables, parser registries, and optional arrays are released;
\item \code{Param_t} is freed late because earlier destructors consult it;
\item reduction buffers and finally \code{Home_t} itself are freed.
\end{enumerate}

The beginning of the implementation shows the alias-first ordering and the
manual walk over raw node blocks:

\lstinputlisting[style=cpp,firstline=40,lastline=137,firstnumber=40]{../src/ParadisFinish.cc}

Preserve this ordering when adding a long-lived subsystem.  A destructor that
needs a control value or lookup table must run before that dependency is
released.

\section{Param\_t and parsing metadata}

\subsection{What Param\_t contains}

Every MPI rank owns one \code{Param_t}.  It combines user controls, derived
constants, function pointers selected during initialization, evolving global
statistics, and output counters.  Most fields are inline scalars or fixed
arrays.  The notable direct heap member is
\code{partialDisloDensity}, allocated according to the selected Burgers-vector
group count and freed before the surrounding \code{Param_t}.

Rank zero allocates \code{ctrlParamList} and \code{dataParamList}.  Each list
owns a growable \code{VarData_t[]} registration array.  A registration's
\code{valList}, however, is a borrowed address, normally into the live
\code{Param_t}; it must not be freed separately.

\begin{longtable}{@{}p{0.26\linewidth}p{0.24\linewidth}p{0.42\linewidth}@{}}
\toprule
Object/member & Storage relationship & Release rule \\
\midrule
\endhead
\code{home->param} & Home-owned object on every rank & Free nested
\code{partialDisloDensity}, then free \code{Param_t} \\
\code{ctrlParamList}, \code{dataParamList} & Home-owned, task-zero registry
objects & Free each \code{varList}, then each list object \\
\code{VarData_t::valList} & Borrowed pointer into registered storage & Never
free through the registry \\
mobility function pointers & Non-owning executable addresses & Re-select on
each rank; no deallocation \\
\bottomrule
\end{longtable}

\subsection{MPI broadcast caveat}

Initialization broadcasts the raw bytes of \code{Param_t}.  Pointer-bearing
parser registries are therefore kept on task zero rather than broadcast.  The
mobility function pointers are selected after the broadcast, and dynamically
sized parameter data is allocated on the receiving rank.  Any new heap-backed
parameter field needs an explicit serialization/allocation policy; adding a
pointer to the structure without one would copy a meaningless process-local
address.

\section{Tag\_t and network identity}

\code{Tag_t} is a small value type containing \code{domainID} and
\code{index}; it owns no memory.  A tag is resolved through one of two sparse
views:

\begin{itemize}
\item a local tag uses \code{home->nodeKeys[tag.index]};
\item a remote tag uses
\code{home->remoteDomainKeys[tag.domainID]->nodeKeys[tag.index]}.
\end{itemize}

Remote resolution may legitimately return null when the rank does not
currently hold that ghost.  Use \code{GetNodeFromTag()} or
\code{GetNodeFromIndex()} instead of duplicating unchecked table traversal.

A tag is more durable than a raw pointer across ordinary local calculations,
but it is not permanent identity across ownership transfer.  Migration and
some initialization/compaction paths assign new tags and build temporary
\code{TagMap_t} value records.  \code{DistributeTagMaps()} communicates and
applies those mappings to neighbor arms, then consumes the map.  Code that
spans migration must participate in remapping or rebuild its keys afterward.

\section{Node\_t: vertex state and endpoint-owned arms}

\subsection{Layout}

\lstinputlisting[style=cpp,firstline=141,lastline=205,firstnumber=141]{../include/Node.h}

A \code{Node_t} is both a dislocation-network vertex and the storage for all
data attached to its incident arms.  Arrays indexed by arm number include the
neighbor tag, per-endpoint force, Burgers vector, glide-plane normal, and
local/remote stress contributions.  Consequently:

\begin{itemize}
\item the persistent graph has no single canonical heap object for a
segment;
\item arm index \(i\) has meaning only while the node's arm arrays keep their
current layout;
\item topology operations must update both endpoints' arm records and preserve
reciprocal neighbor tags, Burgers-vector sign conventions, and compatible
glide planes;
\item inserting, removing, or reallocating arms invalidates saved arm-array
pointers and can shift saved arm indices.
\end{itemize}

Use the topology helpers such as \code{InsertArm()},
\code{ChangeConnection()}, \code{ChangeArmBurg()}, \code{SplitNode()}, and
\code{MergeNode()} rather than editing parallel arrays independently.

\subsection{Two active arm-allocation families}

The codebase contains two distinct ownership paths:

\begin{description}
\item[Legacy production helpers]
\code{AllocNodeArms()}, \code{ReallocNodeArms()}, and \code{FreeNodeArms()}
manage the main simulation arrays: \code{nbrTag}, the Burgers, glide-plane,
and force arrays, plus \code{sigbLoc} and \code{sigbRem}.  A same-degree
allocation reuses capacity and reinitializes entries.  These helpers do not
generally own \code{nbrIndex}, \code{armCoordIndex}, or
\code{armNodeIndex}.

\item[C++ Node\_t methods]
\code{Node_t::Allocate_Arms()}, \code{Free_Arms()}, and \code{Recycle()}
manage all declared per-arm pointers.  The C++ destructor calls
\code{Recycle()}, and copy assignment deep-copies the arrays.
\end{description}

Production pool nodes are raw-\code{calloc}ed and later raw-\code{free}d, so
their constructors and destructors do not run.  The production pool must use
the legacy/manual cleanup path consistently.  Topology backup helpers form an
additional synchronized deep-copy path and contain an in-source warning that
their array lists must stay aligned with \srcfile{src/Util.cc}.

Do not mix these families.  In particular, calling the narrower legacy free
routine after allocating all C++ auxiliary arrays can leak or strand those
auxiliaries.

\subsection{Logical deletion retains capacity}

\lstinputlisting[style=cpp,firstline=1316,lastline=1379,firstnumber=1316]{../src/Util.cc}

\code{FreeNode()} means ``remove this logical native node and recycle its
slot.''  It does not free the object and intentionally does not free its arm
arrays.  The next user of that pool slot may reuse the capacity.  Creation and
unpack paths must therefore initialize every scalar and array entry they rely
on; zero-initialization occurred only when the block was first allocated.

\section{The node pool and its non-owning views}

\subsection{Block allocation}

When the free queue is empty, \code{PopFreeNodeQ()} allocates a raw
\code{NodeBlock_t} header and a zeroed array of
\code{NODE_BLOCK_COUNT} nodes, links the block into
\code{home->nodeBlockQ}, initializes locks, and places all slots on the free
queue.

\lstinputlisting[style=cpp,firstline=18,lastline=60,firstnumber=18]{../src/QueueOps.cc}

Although \code{NodeBlock_t} also has a C++ constructor using \code{new[]}
and a destructor using \code{delete[]}, that is not the main runtime path.
Never destroy a production pool block with the C++ path, or destroy a C++
block with \code{free()}.

\subsection{Pool state machine}

\begin{center}
\fbox{\parbox{0.91\linewidth}{
\centering
raw block slot on \code{freeNodeQ}
\(\longrightarrow\) install as native in \code{nodeKeys}
or as ghost in remote \code{nodeKeys}
\(\longrightarrow\) use through lookup/cell/ghost views
\(\longrightarrow\) detach and recycle
\(\longrightarrow\) same address may represent a different tag
}}
\end{center}

Native creation obtains the lowest recycled tag index when possible; otherwise
it advances the key high-water mark and grows \code{nodeKeys}.  Deletion
places the tag index into a min-heap, nulls the authoritative local key entry,
and returns the node to the free queue.

The address of a pool slot is stable until final teardown, but its
\emph{identity is not}.  A cached \code{Node_t*} can silently become a
different node after recycle and reuse.  Pointer stability is therefore not
object-lifetime stability.

\subsection{Views into the pool}

\begin{longtable}{@{}p{0.26\linewidth}p{0.26\linewidth}p{0.40\linewidth}@{}}
\toprule
View & Meaning & Validity rule \\
\midrule
\endhead
\code{home->nodeKeys} & Sparse authoritative native lookup by local tag index
& Entries can be null; the pointer array may move when grown \\
\code{RemoteDomain_t::nodeKeys} & Sparse lookup of currently materialized
ghosts from one domain & Replaced during ghost communication and destroyed on
decomposition rebuild \\
\code{freeNodeQ}, \code{ghostNodeQ}, \code{nativeNodeQ} & Intrusive queue
heads using \code{Node_t::next} & A node can occupy only one of these queues;
\code{nativeNodeQ} is a compatibility/startup view, not the steady-state
authoritative collection \\
\code{ghostNodeList} & Dense pointer view for threaded ghost traversal & Only
entries below \code{ghostNodeCount} are active; array storage may move on
growth \\
\code{Cell_t::nodeQ} & Intrusive spatial membership using
\code{nextInCell} & Rebuilt by sorting and invalid after cell reset/rebalance \\
\bottomrule
\end{longtable}

\code{PushGhostNodeQ()} updates both the intrusive queue and dense pointer
array.  Its source contract requires one writer and forbids concurrent ghost
traversal during update.  Removed ghosts are not necessarily unlinked from
the whole queue immediately, so the queue alone is not a reliable list of
active ghosts.

\subsection{Ghost batch recycling}

\lstinputlisting[style=cpp,firstline=146,lastline=169,firstnumber=146]{../src/QueueOps.cc}

Batch recycling splices ghost slots onto the free queue and resets only the
active ghost count and queue heads.  It does not free arms or clear all old
pointer values.  Consumers must honor counts and phase boundaries, not merely
test whether an old pointer value is non-null.

\section{Segments are usually derived views}

The persistent network edge is represented twice, once in each endpoint's
arm arrays.  \code{Segment_t} materializes a convenient record containing
borrowed \code{Node_t*} endpoints and copied positions, force values, Burgers
vector, glide plane, cell, and native flags.  Its destructor destroys its
lock when applicable; it never deletes either endpoint.

\code{SegmentPair_t} borrows two \code{Segment_t*} records and owns no nested
memory.  Thus the common hierarchy is:

\begin{center}
\fbox{\parbox{0.88\linewidth}{
caller owns segment-array storage
\(\rightarrow\) each \code{Segment_t} borrows pooled nodes
\(\rightarrow\) a pair list owns or borrows pair-array storage
\(\rightarrow\) each pair borrows segment records
}}
\end{center}

The allocator for the outer list determines its release rule.  The C++
\code{Segment_List_Create()} and \code{Segment_Append()} family grows arrays
with \code{new Segment_t[]} and releases replaced arrays with
\code{delete[]}; callers must ultimately \code{delete[]} the returned list.
Some force kernels instead build C-allocated scratch segment arrays and free
them before returning.  Follow the producing API rather than assuming all
\code{Segment_t*} pointers have the same allocator.

No segment view may outlive its endpoint population.  Rebuild segment and
pair lists after topology, migration, or ghost refresh.

\section{Cells, collision grids, and inclusions}

\subsection{Persistent Cell\_t storage}

\lstinputlisting[style=cpp,firstline=11,lastline=41,firstnumber=11]{../include/Cell.h}

\code{home->cellTable} is a hash table that owns all allocated
\code{Cell_t} records.  \code{AddCellToTable()} uses \code{calloc()} and
links a new cell into a bucket.  Cells intentionally remain in the hash table
for the whole run to avoid repeated structural allocation.

A cell owns its \code{nbrList} and \code{domains} integer arrays.  It does not
own the nodes reachable through \code{nodeQ}; those are pooled nodes linked by
\code{nextInCell}.  With Eshelby support, \code{inclusionQ} is an integer
index into the flat Home-owned inclusion array rather than a pointer.

On load rebalance, \code{DLBfreeOld()} retains cell structs and static
neighbor lists but clears active flags and queues, frees dynamic domain lists
and \code{cellList}, and releases remote-domain metadata.  The new
decomposition rebuilds these views.  At final shutdown,
\code{FreeCellTable()} releases cell-owned arrays, locks, and cell records.

\subsection{Collision cell2}

The collision overlay is Home-owned scratch.  \code{home->cell2} stores queue
heads, and \code{cell2QentArray} stores entries containing borrowed node
pointers plus integer next links.  Sorting rebuilds it for collision use;
\code{FreeNode()} removes nodes from this view before recycling them.  Neither
array owns a node.

\subsection{Eshelby inclusion storage}

\code{EInclusion_t} is a pointer-free value record.  The Home-owned inclusion
array, \code{eshelbyInclusions}, is contiguous and reallocatable.  Locally
owned inclusions come first and ghost inclusions follow them.  Migration and
ghost exchange can compact or reallocate the array, so code must not retain
an \code{EInclusion_t*} across those phases.  Cell queues use integer indices
and are rebuilt after ownership changes.

\code{segpartIntersectList} is a growable, phase-scoped list keyed by copied
segment tags and inclusion indices.  Its producer serializes appends with the
associated lock and clears/frees the list after use.  Saved pointers into the
list are invalid after a growth operation.

\section{MPI domains, ghosts, and migration}

\subsection{RemoteDomain\_t ownership}

\lstinputlisting[style=cpp,firstline=16,lastline=31,firstnumber=16]{../include/RemoteDomain.h}

\code{home->remoteDomainKeys} owns a domain-ID-indexed table of remote-domain
records, while \code{home->remoteDomains} is a dense list of active IDs with
primary neighbors before secondary-only neighbors.  Each remote record owns:

\begin{itemize}
\item its \code{expCells} array when present;
\item the storage for its sparse \code{nodeKeys} pointer table, but not the
pooled ghost nodes referenced by its entries;
\item \code{inBuf} and \code{outBuf} only for the protocol phase that
allocated or received them.
\end{itemize}

Persistent Home-level MPI request and status arrays are allocated when remote
domains are initialized.  Protocol-specific buffers must remain allocated
until the matching nonblocking requests complete.

The fields \code{inBufLen} and \code{outBufLen} are protocol-polymorphic:
several ghost, velocity, and remesh paths treat them as bytes, while
secondary-ghost phases use element counts.  Interpret a length only within
the protocol that set it.

\subsection{End-of-cycle ordering}

\lstinputlisting[style=cpp,firstline=525,lastline=541,firstnumber=525]{../src/ParadisStep.cc}

This sequence is a lifetime contract:

\begin{enumerate}
\item \code{Migrate()} transfers natives that crossed a domain boundary.
\item \code{RecycleGhostNodes()} makes old ghost pool slots reusable.
\item \code{SortNativeNodes()} clears and rebuilds cell membership from the
authoritative native key table.
\item \code{CommSendGhosts()} replaces remote lookup tables and materializes
the new ghost population.
\end{enumerate}

Between recycle and replacement, old remote lookup storage or old
\code{ghostNodeList} slots may still contain pointer values to reusable pool
slots.  They are not valid ghosts.  Do not resolve remote tags or traverse
stale entries in that interval.

\subsection{Migration transfers logical ownership, not storage blocks}

Migration performs a value transfer:

\begin{enumerate}
\item the source selects outgoing native nodes and serializes their scalar and
arm data into caller-owned send buffers;
\item after packing a node, the source calls \code{FreeNode()}, making its
local pool slot and tag index reusable while retaining arm capacity;
\item the destination obtains a slot from its own local pool, assigns a new
local tag, allocates or reuses arms, and unpacks values;
\item the destination records \code{oldTag -> newTag};
\item \code{DistributeTagMaps()} repairs neighbor tags across ranks and frees
the temporary mapping storage;
\item send/receive arrays are released only after communication completes.
\end{enumerate}

No \code{Node_t} allocation crosses MPI.  The source and destination each own
their own storage, and only serialized values cross the wire.

\subsection{Operation ownership for cross-domain segments}

\code{DomainOwnsSeg()} gives same-domain segments to that domain.  For a
cross-domain segment it alternates the responsible rank by even/odd cycle;
remesh responsibility is the inverse of collision/separation responsibility.
This creates deterministic, balanced mutation authority.  It has no bearing
on node-block, arm-array, or communication-buffer deallocation.

\subsection{Topology operation buffers}

\code{home->opListBuf} is a persistent, growable byte buffer.  Initialization
allocates an initial capacity, \code{ClearOpList()} resets used length while
retaining capacity, and \code{FreeOpList()} releases it at shutdown.
Operation records are packed value snapshots, not owners of live pointers.

The remesh communication path has a special alias: one send-buffer entry can
refer directly to \code{home->opListBuf}, while other entries are allocated
copies.  The communication cleanup must not free the aliased persistent
entry.  Receive-buffer ownership is transferred to \code{FixRemesh()}, which
consumes and frees it.

\subsection{Mirror domains}

Only task zero creates \code{MirrorDomain_t} snapshots for plotting/output.
A mirror owns its key-array storage, remote-arm coordinate arrays, and optional
copied inclusion array.  Its node-key entries point to nodes obtained from the
same common pool; \code{MirrorDomain_Free()} releases mirror-owned arrays and
returns those nodes to the free queue.  Mirror nodes are output snapshots, not
primary ghosts, and their pointers are invalid after mirror cleanup.

\section{Other owned subsystems}

\subsection{Initialization staging: InData\_t}

\code{InData_t} exists only during startup.  It owns a block-read node array
and those nodes' arm arrays, compatibility Burgers arrays for older inputs,
and the input decomposition.  Its \code{param} field borrows the live
\code{Param_t}.  \code{FreeInitArrays()} deep-frees the staging allocations,
then initialization frees the container.  The steady-state
\code{home->decomp} is independently constructed/broadcast and is not an
alias of the freed staging decomposition.

\subsection{Generic decomposition pointer}

\code{home->decomp} is an owned \code{void*} whose concrete layout depends on
the selected decomposition type.  Recursive sectioning owns nested boundary
arrays; recursive bisection owns a tree.  Always release it through
\code{FreeDecomp(home, home->decomp)}.  A plain \code{free()} cannot traverse
either representation correctly.

\subsection{FMM hierarchy}

When enabled, Home owns a \code{FMLayer_t[]} hierarchy.  Each layer owns
communication-domain arrays, a cell-ID list, and hash entries.  Each
\code{FMCell_t} owns domain-intersection and multipole/expansion coefficient
arrays; its next/previous fields are only hash links.  Initial setup creates
the static hierarchy, while load-balance refreshes recompute dynamic
intersection data.  Treat \code{FMInit()} and \code{FMFree()} as the
constructor/destructor boundary rather than freeing nested fields from force
callers.

\subsection{Subcycling state}

\code{Subcycling_t} is deliberately opaque outside
\srcfile{src/SubcyclingIntegrator.cc}.  It is a true C++ owner of STL maps and
vectors and is created with \code{new} by \code{StartSubcycling()} and
destroyed with \code{delete} by \code{FinishSubcycling()}.  Its contract limits
it to one outer cycle and ends it before topology or migration can invalidate
its node indices, tag-pair keys, and borrowed node observations.

\subsection{Material and solver state}

Home embeds \code{BurgInfo_t}, whose dynamically generated Burgers-vector,
plane, indexing, and splinterability arrays are released individually during
shutdown.  Optional anisotropic, spectral, ARKode, Eshelby-FMM, and related
subsystems likewise have feature-specific cleanup.  New conditional storage
should have symmetric initialization and release under the same preprocessor
guard.

\section{Allocation and release matrix}

\small
\begin{longtable}{@{}p{0.20\linewidth}p{0.20\linewidth}p{0.24\linewidth}p{0.28\linewidth}@{}}
\toprule
Storage & Allocator/creator & Storage owner & Release or reuse path \\
\midrule
\endhead
\code{Home_t} & \code{calloc} in \code{InitHome()} & main simulation driver &
\code{ReleaseMemory()} then \code{free(home)} \\
\code{Param_t} & \code{calloc} during \code{Initialize()} & \code{Home_t} &
free nested density array, then \code{free(param)} \\
parser registries & \code{calloc}; \code{BindVar()} grows \code{varList} &
\code{Home_t}, task zero & free \code{varList} and registry, never
\code{valList} targets \\
node-block header and slots & \code{malloc} header plus \code{calloc} slots in
\code{PopFreeNodeQ()} & \code{home->nodeBlockQ} & manual per-slot arm/lock
cleanup, then \code{free} slots and header \\
runtime node arms & \code{AllocNodeArms()} / \code{ReallocNodeArms()} & owning
\code{Node_t} slot & normally retained on \code{FreeNode()};
\code{FreeNodeArms()} or shutdown \\
native key array and recycled heap & \code{realloc}, heap helpers &
\code{Home_t} & final cleanup; entries never own nodes \\
ghost traversal array & \code{realloc} in \code{PushGhostNodeQ()} &
\code{Home_t} & capacity retained across cycles, freed at shutdown \\
\code{Cell_t} records & \code{calloc} in \code{AddCellToTable()} & cell hash
table & \code{FreeCellTable()} \\
cell dynamic arrays & cell/decomposition initialization & owning cell or Home
epoch state & \code{DLBfreeOld()} as appropriate, final cell cleanup \\
\code{RemoteDomain_t} and arrays & remote-domain initialization &
\code{remoteDomainKeys} table & \code{DLBfreeOld()}; never free ghost nodes
through key entries \\
protocol buffers & per communication routine & allocating protocol or
transferred consumer & after matching MPI completion and unpack/consume \\
operation-list buffer & \code{InitOpList()}, grow on pack & \code{Home_t} &
clear retains capacity; \code{FreeOpList()} releases \\
segment list from C++ list API & \code{new Segment_t[]} & caller &
\code{delete[]} after all segment pairs are done \\
mirror domain arrays & mirror receive/unpack & task-zero mirror table &
\code{MirrorDomain_Free()} and surrounding table cleanup \\
\code{InData_t} staging arrays & restart/input readers & \code{InData_t} &
\code{FreeInitArrays()}, then free container \\
decomposition & type-specific allocators & \code{Home_t} & \code{FreeDecomp()} \\
FMM layers/cells/coefficients & \code{FMInit()} and helpers & \code{Home_t} /
nested FMM records & \code{FMFree()} \\
subcycling state & \code{new Subcycling_t} & \code{Home_t} for one cycle &
\code{delete} in \code{FinishSubcycling()} \\
inclusion array & input, ghost, and migration paths using \code{realloc} &
\code{Home_t} & final free; indices rebuilt after movement \\
\bottomrule
\end{longtable}
\normalsize

\section{Pointer and index invalidation rules}

\begin{longtable}{@{}p{0.29\linewidth}p{0.29\linewidth}p{0.34\linewidth}@{}}
\toprule
Saved value & Invalidated by & Safe replacement \\
\midrule
\endhead
\code{Node_t*} & node removal/recycle, ghost refresh, migration, mirror cleanup
& save a tag only within a stable tag epoch and re-resolve; rebuild after
migration \\
pointer to \code{nodeKeys} entry/storage & key-array growth via
\code{realloc}, DLB rebuild & re-read \code{home->nodeKeys} and bounds \\
remote-node pointer or key-table entry & ghost recycle/replacement, DLB &
resolve only during a valid ghost epoch \\
arm-array pointer or arm index & insert/delete/reallocation of arms & locate
the neighbor arm again by tag \\
cell queue traversal state & sorting, movement, DLB, node deletion & restart
from rebuilt cell queues \\
\code{Segment_t*} or \code{SegmentPair_t*} & segment-list growth/destruction,
topology, endpoint recycle & rebuild views and pairs \\
\code{EInclusion_t*} & inclusion-array growth or compaction & store stable
semantic ID where appropriate, then locate again; rebuild cell indices \\
pointer into tag map, op list, or intersection list & any append that grows
the backing array, clear/free & use copied value/offset only under the API's
documented phase \\
MPI send/receive buffer & completion cleanup or ownership transfer & do not
access after the consumer frees it; never move/free before request completion \\
subcycling cache/index & topology, tag changes, node-key capacity/layout
changes & keep state inside the prescribed cycle and rebuild next cycle \\
\bottomrule
\end{longtable}

\section{Memory API and diagnostics}

Translation units including \srcfile{include/Home.h} replace C
\code{malloc()}, \code{calloc()}, \code{realloc()}, and \code{free()} calls
with \code{ParadisMalloc()}, \code{ParadisCalloc()},
\code{ParadisRealloc()}, and \code{ParadisFree()}, including source location
metadata.

\lstinputlisting[style=cpp,firstline=130,lastline=149,firstnumber=130]{../include/Home.h}

Without \code{DEBUG_MEM}, these wrappers delegate to the system allocator and
fail fast on allocation errors.  With \code{DEBUG_MEM}, they track blocks and
guard regions; the timestep driver calls \code{ParadisMemCheck()} at selected
phase boundaries, and termination reports/cleans the tracking layer after
\code{ParadisFinish()}.

C++ \code{new}/\code{delete} do not pass through these macros.  Pair
allocation families exactly:

\begin{itemize}
\item \code{malloc/calloc/realloc} with \code{free};
\item \code{new T} with \code{delete};
\item \code{new T[n]} with \code{delete[]};
\item subsystem creators with their documented subsystem destructor.
\end{itemize}

\section{Current implementation sharp edges}

The following are implementation observations, not recommended ownership
patterns.  They matter when modifying nearby code.

\begin{enumerate}
\item \code{GetNewNativeNode()} resets only selected fields and does not
release retained arm capacity.  Current topology callers explicitly clear
arms and set native state.  Every new caller must fully initialize its slot.

\item All free, ghost, and compatibility native queues share
\code{Node_t::next}.  A node cannot be linked in two of them at once.
\code{PushFreeNodeQ()} does not update \code{lastFreeNode} when pushing into
an empty queue, while ghost-batch splicing relies on that tail; changes to pool
transition ordering should audit this invariant.

\item Primary remote-domain records are allocated with \code{malloc()} and
only a subset of fields is initialized immediately.  Existing control flow
uses count fields before pointer fields.  New users should not assume every
pointer is null merely because the Home root was zeroed.

\item Not every communication cleanup path nulls freed
\code{RemoteDomain_t::inBuf/outBuf} fields or resets lengths.  The values are
valid only inside the protocol's documented phase, even if a stale non-null
bit pattern remains afterward.

\item Normal final cleanup assumes transient \code{tagMap} and subcycling
state have already been consumed and that mirror-table lifecycle completed.
The persistent MPI request/status arrays allocated during remote-domain
initialization are not visibly released by \code{ReleaseMemory()} in this
revision.  Leak-clean early-exit work should audit these paths explicitly.

\item \code{ReleaseMemory()} relies on manual lists of conditional fields.
When adding a pointer to \code{Home_t}, \code{Param_t}, \code{Node_t},
\code{Cell_t}, or an optional subsystem, add allocation, reset/invalidation,
and normal plus early-exit cleanup together.
\end{enumerate}

\section{Extension checklist}

Before adding or retaining a pointer, buffer, or index, answer all of the
following:

\begin{enumerate}
\item Who owns the storage, and which exact function releases it?
\item Is the pointer an object, an intrusive link, a lookup entry, or a
phase-scoped borrowed view?
\item Which event invalidates it: \code{realloc}, topology, node recycle,
migration, ghost refresh, cell sort, DLB, or shutdown?
\item Does the allocation use the C wrappers, scalar \code{new}, array
\code{new[]}, or a subsystem allocator, and is the release family identical?
\item If the field is in \code{Param_t}, how is it reconstructed after raw
MPI broadcast?
\item If it is an MPI buffer, who owns it before and after each nonblocking
request, and in what units is its length stored?
\item If it refers to a node or segment, can a copied \code{Tag_t} plus later
resolution replace a long-lived raw pointer?
\item If an array grows, are all interior pointers and saved indices either
rebuilt or prohibited?
\item Is initialization complete for a recycled/raw-allocated object rather
than relying on a constructor or old zeroes?
\item Are feature guards, locks, failure paths, DLB rebuild, and final cleanup
symmetric?
\end{enumerate}

The safest steady-state node traversal is the sparse local key table:

\begin{lstlisting}[style=cpp,numbers=none]
for (int i = 0; i < home->newNodeKeyPtr; ++i) {
    Node_t *node = home->nodeKeys[i];
    if (node == NULL) continue;
    /* Use node only within the current topology/migration/ghost epoch. */
}
\end{lstlisting}

Across a potentially invalidating phase, discard the pointer and resolve
again from a valid, remapped tag.  Across migration itself, rebuild the
relation after tag maps have been distributed.

\section{Summary of invariants}

\begin{itemize}
\item \code{Home_t} is the practical per-rank storage root, but nested
subsystems still require specialized teardown.
\item \code{nodeBlockQ}, not lookup tables or queues, owns production node
objects.
\item A node owns its active or retained arm allocations; logical deletion
usually recycles rather than deallocates them.
\item \code{nodeKeys}, ghost lists, remote key tables, cell queues, segment
records, and segment pairs are views with narrower validity windows.
\item Ghost memory is local even when logical ownership is remote.
\item MPI operation authority and segment-computation ownership are unrelated
to heap ownership.
\item Cells persist, while decomposition-epoch lists and associations are
cleared and rebuilt.
\item Migration serializes values into independently owned destination
storage, changes tags, repairs references, and invalidates prior relations.
\item Raw allocation of class types bypasses constructors/destructors, so the
manual allocator family and shutdown order are part of the correctness model.
\end{itemize}

\end{document}
